\documentclass[10pt,a4paper]{amsart}
\usepackage[margin=19mm]{geometry}
\usepackage{amsmath,amssymb,mathtools}
\usepackage[scr=rsfs,cal=euler]{mathalfa}
\usepackage{xfrac}
\usepackage[unicode,hidelinks,bookmarks=false]{hyperref}
\newcommand{\ie}{i.e.\ }
\newcommand{\cf}{cf.\ }
\newcommand{\resp}{respectively\ }
\newcommand{\Subsection}{\subsection}
\providecommand{\nobreakdash}{\nobreak\hbox{-}}
\setlength{\emergencystretch}{2em}
\multlinegap0pt
\usepackage{bm}
\usepackage{bbm}
\usepackage{dsfont}
\usepackage{xspace}
\usepackage{cancel}
\usepackage{mathtools}
\usepackage{amsthm}
\makeatletter
\@ifundefined{theorem}{\newtheorem{theorem}{Theorem}}{}
\@ifundefined{lemma}{\newtheorem{lemma}[theorem]{Lemma}}{}
\makeatother
\title[Quasilinear elliptic problems with singular nonlinearities i]{Quasilinear elliptic problems with singular nonlinearities in half-spaces}
\author{Phuong Le}
\date{Source: arXiv:2409.19557v2 (CC BY 4.0)}
\begin{document}
\maketitle

\noindent\emph{Source and licence.} Quasilinear elliptic problems with singular nonlinearities in half-spaces. Version arXiv:2409.19557v2 (CC BY 4.0), distributed under the Creative Commons Attribution 4.0 International licence, as recorded in the arXiv licence metadata for this exact version and shown on the abstract page https://arxiv.org/abs/2409.19557v2. Full rights notice: licenses/arxiv-2409.19557-CC-BY-4.0.txt.\par\smallskip

\noindent\textit{Editorial scope.} Counted unit: the Lemma whose source label is \texttt{lem:lower\_bound} in the article (source0.tex lines 613--619, source label \texttt{lem:lower\_bound}), delivered together with the article's own complete original proof of it (source0.tex lines 621--696, one \texttt{proof} environment of 76 lines). No mathematics is written, repaired or re-derived: statement and proof are the article's own text line for line. Required context retained uncounted: main (lines 69--75); p\_ge (lines 433--436). Every definition, display and label of those lines is retained unchanged. Omitted from this package and not redistributed: 1--68, 76--432, 437--612, 620--620, 697--1688. Nothing inside a retained excerpt is omitted or altered: every retained body is delivered exactly as the article's lines are written, so no omission receipt of any kind accompanies this package. Citations: the retained excerpts carry no citation command at all, so no bibliography entry of the article is transcribed and no reference is redistributed. Reference adaptation: none was needed and none was made; every reference inside the delivered unit resolves inside it, so no unresolved reference or citation is produced. Printed slips kept literal: no printed word, symbol, hypothesis or number of the counted statement or of its proof was altered. Layout and class: the article's own class is \texttt{amsart} with packages including hyperref, cite, enumitem; the wrapper is the lane's pinned \texttt{amsart} editorial wrapper at 10pt on A4, which restates the theorem-like environments the retained text needs and adds the article's own macro definitions that the retained text needs (none), with extra wrapper packages (bm, bbm, dsfont, xspace, cancel, mathtools). The delivered numbering is the unit's own. Assets: no retained span contains a figure environment, a graphics-inclusion command, a PDF-page inclusion, a table or a TikZ picture; no image file is copied into this package, the package declares no asset dependency and no image file is redistributed with it. Licence: version arXiv:2409.19557v2 of this article is distributed under the Creative Commons Attribution 4.0 International licence, as recorded in the arXiv licence metadata for this exact version and shown on the abstract page https://arxiv.org/abs/2409.19557v2. A full rights notice is delivered at licenses/arxiv-2409.19557-CC-BY-4.0.txt. Characters: every retained excerpt is pure ASCII, so the delivered engine, XeLaTeX with its default Latin Modern text font, reports no missing character.
	\begin{equation}\label{main}
		\begin{cases}
			-\Delta_p u = f(u) & \text{ in } \mathbb{R}^N_+,\\
			u>0 & \text{ in } \mathbb{R}^N_+,\\
			u=0 & \text{ on } \partial\mathbb{R}^N_+,
		\end{cases}
	\end{equation}

	\begin{align}
		(|\xi|^{p-2}\xi - |\xi'|^{p-2}\xi', \xi -\xi') &\ge C_1 (|\xi| + |\xi'|)^{p-2} |\xi -\xi'|^2, \label{p_ge}\\
		\left||\xi|^{p-2}\xi - |\xi'|^{p-2}\xi'\right| &\le C_2 (|\xi| + |\xi'|)^{p-2} |\xi -\xi'|, \label{p_le}
	\end{align}

	\begin{lemma}\label{lem:lower_bound}
		Let $f:(0,+\infty)\to\mathbb{R}$ be a locally Lipschitz continuous function such that $f(t) > \frac{c_0}{t^\gamma}$ for all $0<t<t_0$, where $c_0,t_0>0$, $\gamma\ge0$. Let $u \in C^{1,\alpha}_{\rm loc}(\mathbb{R}^N_+) \cap C(\overline{\mathbb{R}^N_+})$ be a solution to \eqref{main}. Then
		\[
		u(x) \ge \min\{C x_N^\frac{p}{\gamma+p-1}, t_0\} \quad\text{ in } \mathbb{R}^N_+.
		\]
		for some constant $C>0$.
	\end{lemma}

	\begin{proof}
		Let $\lambda_1>0$ and $\phi_1 \in C^1(\overline{B_1})$ be the first eigenvalue and a corresponding positive eigenfunction of the $p$-Laplacian in $B_1$, namely,
		\[
		\begin{cases}
			-\Delta_p \phi_1 = \lambda_1 \phi_1^{p-1} & \text{ in } B_1,\\
			\phi_1 > 0 & \text{ in } B_1,\\
			\phi_1 = 0 & \text{ on } \partial B_1.
		\end{cases}
		\]
		
		Setting
		\[
		w = s \phi_1^\frac{p}{\gamma+p-1},
		\]
		where $s>0$ will be chosen later. Direct calculation yields that in the weak sense
		\[
		-\Delta_p w = \frac{\alpha(x)}{w^\gamma} \quad \text{ in } B_1,
		\]
		where
		\[
		\alpha(x) := s^{\gamma+p-1} \left(\frac{p}{\gamma+p-1}\right)^{p-1} \left[\frac{(\gamma-1)(p-1)}{\gamma+p-1}|\nabla\phi_1(x)|^p + \lambda_1\phi_1(x)^p\right].
		\]
		
		Now we fix $s>0$ such that $\sup_{x\in B_1}\alpha(x) \le c_0$ and hence
		\[
		-\Delta_p w \le \frac{c_0}{w^\gamma} \text{ in } B_1.
		\]
		Let $R_0>0$ be such that $R_0^\frac{p}{\gamma+p-1} w(0) = t_0$.
		
		For any $0<R\le R_0$ and $x_0 = (x_0', x_{0,N})\in\mathbb{R}^N_+$ with $x_{0,N} \ge R + \varepsilon$, where $\varepsilon$ is sufficiently small, we set
		\[
		w_{x_0,R}(x) := R^\frac{p}{\gamma+p-1}w\left(\frac{x-x_0}{R}\right) \quad\text{ in } B_R(x_0).
		\]
		Then
		\[
		w_{x_0,R} \le t_0 \quad\text{ and }\quad -\Delta_p w_{x_0,R} \le \frac{c_0}{w_{x_0,R}^\gamma} \quad\text{ in } B_R(x_0).
		\]
		
		On the other hand, since $w_{x_0,R}=0 < u$ on $\partial B_R(x_0)$, we can use $(w_{x_0,R}-u)^+ \chi_{B_R(x_0)}$ as a test function in
		\[
		-\Delta_p u = f(u)
		\quad\text{ and }\quad
		-\Delta_p w_{x_0,R} \le \frac{c_0}{w_{x_0,R}^\gamma}
		\]
		to obtain
		\begin{align*}
			&\int_{B_R(x_0)} (|\nabla w_{x_0,R}|^{p-2}\nabla w_{x_0,R} - |\nabla u|^{p-2}\nabla u, \nabla (w_{x_0,R}-u)^+) \\
			&\le \int_{B_R(x_0)} \left(\frac{c_0}{w_{x_0,R}^\gamma} - f(u)\right) (w_{x_0,R}-u)^+.
		\end{align*}
		
		In $B_R(x_0)\cap\{w_{x_0,R} > u\}$ we have $f(u) \ge \frac{c_0}{u^\gamma}$. Hence
		\begin{align*}
			&\int_{B_R(x_0)} (|\nabla w_{x_0,R}|^{p-2}\nabla w_{x_0,R} - |\nabla u|^{p-2}\nabla u, \nabla (w_{x_0,R}-u)^+) \\
			&\le \int_{B_R(x_0)} \left(\frac{c_0}{w_{x_0,R}^\gamma} - \frac{c_0}{u^\gamma}\right) (w_{x_0,R}-u)^+ \le 0.
		\end{align*}
		
		By \eqref{p_ge}, this implies
		\[
		\int_{B_R(x_0)} (|\nabla w_{x_0,R}| + |\nabla u|)^{p-2} |\nabla (w_{x_0,R}-u)^+|^2 \le 0.
		\]
		Hence $u \ge w_{x_0,R}$ in $B_R(x_0)$ with $x_{0,N} \ge R + \varepsilon$. Since $\varepsilon>0$ is arbitrary, we deduce
		\[
		u \ge w_{x_0,R} \text{ in } B_R(x_0) \text{ for all } 0<R\le R_0 \text{ and } x_0 \in \mathbb{R}^N_+ \text{ with } x_{0,N} \ge R.
		\]
		
		In particular, if $x_{0,N} = R < R_0$, then
		\[
		u(x_0) \ge w_{x_0,R}(x_0) = w(0)R^\frac{p}{\gamma+p-1} = w(0) x_{0,N}^\frac{p}{\gamma+p-1}.
		\]
		If $x_{0,N} \ge R = R_0$, then
		\[
		u(x_0) \ge w_{x_0,R}(x_0) = w(0)R_0^\frac{p}{\gamma+p-1} = t_0.
		\]
		
		The conclusion follows from the fact that $x_0$ is chosen arbitrarily in $\mathbb{R}^N_+$.
	\end{proof}

\end{document}
